\begin{table}[htbp]
\centering
\caption{Standardized Effect Sizes}
\label{tab:sde}
\begin{adjustbox}{max width=\textwidth}
\begin{tabular}{lcccccc}
\hline\hline
Outcome & $\hat{\beta}$ & SE & SD($Y$) & SDE & SE(SDE) & Classification \\
\hline
All fatal crashes & 0.2052 & 0.2743 & 1.5071 & 0.1361 & 0.1820 & Moderate positive \\
Phone-distracted crashes & 0.0089 & 0.0425 & 0.4580 & 0.0194 & 0.0928 & Small positive \\
Any distraction crashes & 0.0318 & 0.0847 & 0.9215 & 0.0345 & 0.0919 & Small positive \\
\hline\hline
\end{tabular}
\end{adjustbox}
\begin{tablenotes}
\item \textit{Notes:} \textbf{Country:} United States. \textbf{Research question:} Do state-level handheld cellphone bans reduce fatal traffic crashes near state borders where the policy discontinuously changes? \textbf{Policy mechanism:} Handheld cellphone bans prohibit drivers from holding a phone while operating a vehicle, aiming to reduce distracted driving by shifting phone use to hands-free modes or eliminating it; enforcement typically involves primary-offense traffic stops. \textbf{Outcome definition:} Monthly counts of fatal motor vehicle crashes (FARS), phone-distracted fatal crashes (FARS distraction codes 5, 6, 15), and any-distraction fatal crashes at the county level within 30km of treated state borders. \textbf{Treatment:} Binary; state-level adoption of handheld cellphone ban between 2017--2021. \textbf{Data:} NHTSA Fatality Analysis Reporting System, 2015--2022, county-month panel, 5,875 observations across 234 counties in 8 border pairs. \textbf{Method:} Difference-in-discontinuities with border-pair and year-month fixed effects; standard errors clustered at state-county level; robustness via wild cluster bootstrap, donut RDD, and distance polynomial controls. \textbf{Sample:} Fatal crashes within 30km of state borders where one state adopted a handheld cellphone ban and the neighboring state did not; restricted to continental U.S. with valid geocoded coordinates. SDE $= \hat{\beta} / \text{SD}(Y)$ where SD($Y$) is the pre-treatment standard deviation. Classification refers to magnitude, not statistical significance: Large ($|$SDE$| > 0.15$), Moderate ($0.05$--$0.15$), Small ($0.005$--$0.05$), Null ($< 0.005$).
\end{tablenotes}
\end{table}

